\section{Summary and conclusions}\label{sec:conclusions}
%%%%%%%%%%%%%%

The CTAO will be the next generation of ground-based $\gamma$-ray Cherenkov detectors, offering a unique opportunity to study high-energy phenomena in the TeV range.  In addition to targeted observations, CTAO is expected to carry out large-area sky surveys, making it particularly well-suited for anisotropy-based techniques such as auto-correlation and cross-correlation analyses.
In this work, we have investigated the sensitivity of future CTAO data to constrain the properties of unresolved blazars and DM signals using these methods.

We found, in particular, that the auto-correlation APS method is better suited for probing the signal from unresolved blazars, with its sensitivity strongly dependent on the effective {point-}source detection threshold. On the other hand, the cross-correlation APS technique, when used in combination with galaxy catalogs, offers very promising sensitivity to potential DM annihilation or decay signal, reaching values around $10^{-23}\,\mathrm{cm^{-3}\,s^{-1}}$ for annihilation cross-section $\langle \sigma v \rangle$ and $10^{27}\,\mathrm{s}$ for decay.
These sensitivities are comparable to and complementary with those obtained from existing strategies, such as observations of dwarf spheroidal galaxies or galaxy clusters.


In the future, these results can be further improved with upcoming, deeper, and more densely sampled galaxy catalogs, such as those from the European Space Agency's \textit{Euclid} satellite mission \cite{2024arXiv240513491E}, as well as by incorporating cross-correlations with other large-scale-structure tracers, such as weak lensing.

